\scanpage{16}

Transformons l'équation
\[
x^2=y^3+\unsure{\dot a|a}xy+by^2+cx+dy+e
\]
par $s$, et en rappelant que $sx=-x$, $sy=y$, on trouve
\[
x^2=y^3\unsure{-|+}axy+by^2-cx+dy+e
\]
et en comparant on trouve $a=c=0$, donc l'équation ici se met sous la forme
\[
\boxed{x^{\unsure{2|3}}=y^{\unsure{3|2}}+ay^2+by+c}
\]
avec comme condition $\delta(a,b,c)\neq0$, qui ici ne fait qu'exprimer que les solutions de $P(y)=y^3+ay^2+by+c$ sont distinctes en tt pt de $S$.

\margin{\unsure{quotient}}
Géométriquement, \deleted{\unsure{le schéma}} l'équation s'interprète ainsi : le quotient $C/\mathbb Z_2$ est un schéma de \unsure{Brauer-Sévéri} \inserted{$P$} sur $S$, muni d'une section \unsure{marquée} (correspondant à la section nulle de $C$ sur $S$), donc \deleted{s'identifie} \deleted{est} \deleted{isomorphe}\inserted{canoniquement} à un espace principal homogène sous un \deleted{\unsure{faisceau inverse}} \deleted{fibré \unsure{qui est} loc. libre} \unsure{fibré vectoriel} \deleted{\unsure{de rang un}}\inserted{inversible} [il correspond à $L^{\otimes2}$]

$C$ est ramifié \deleted{et} au dessus d'un diviseur relatif \inserted{de $P/S$} étale sur $S$, qui se décompose en la section à l'infini, et en un diviseur relatif \inserted{de $P/S$} étale sur $S$, \deleted{\unsure{dans les fibres}} de \deleted{\unsure{$E$}} $\underline V(L^{\otimes2})/S$ étale sur $S$, de degré relatif $3$.
\mspage{15}
